\section{A Feasibility-First Workflow for Nonlinear Updates}\label{app:workflow}
Fix a profile space, an acquisition menu, a confidence level $1-\alpha>1/2$,
and a target half-width $e$. The workflow separates two certificates:
a dual lower bound $\delta_L$ on the identification floor and a feasible
approximation $v\in\Hh_\A$ with residual at most $\varepsilon$.
Only the second certificate permits an accuracy guarantee.

\paragraph{Rejecting an insufficient cap.}
If $\delta_L>e$, no interval with half-width at most $e$ can have uniform
coverage $1-\alpha$, for any number of records. To see this on a finite
space, choose a dual signed measure $h$ with $\|h\|_1=1$,
$h\perp\Hh_\A$, and $\sum_xh(x)d(x)>e$.
Define
\[
 P_\pm(x,y)=\frac{|h(x)|\pm yh(x)}2.
\]
Every interview cell has the same joint probability with $Y$ under the
two laws because its sum of $h$ is zero. This also gives identical
experiments under sequential design. Their target values are opposite
and more than $2e$ apart. Under the common observation law an interval
of that width cannot cover both, so the two coverage probabilities sum
to at most one. This rules out uniform coverage greater than one half.
The conclusion assumes arbitrary joint laws on the stated profile space;
a substantively justified restriction on those laws changes the problem.

\paragraph{Passing with a certified approximation.}
Suppose an observable score $H$ has conditional mean $Yv(X)$ and satisfies
$|H|\le C$. Conditional on all model fitting and calibration, fresh iid
records satisfy
\[
 \Pr_P\!\left\{|\overline H-\Delta(P)|
 \le\varepsilon+C\sqrt{2\log(2/\alpha)/n}\right\}\ge1-\alpha.
\]
This is Hoeffding's inequality plus
$|\E_P[Y(v-d)]|\le\varepsilon$. It is valid for every population on the
certified space, not merely the reference population used to fit a score.
For $\varepsilon<e$, the sufficient count in Section~2 follows by inversion.
If the lower and upper certificates straddle $e$, they have not resolved
feasibility: improve the approximation or compare a larger cap.

On a binary cube, write $v=\sum_{|S|\le b}\gamma_S\chi_S$ and
$C=\sum_S|\gamma_S|$. Query term $S$ with probability $|\gamma_S|/C$ and
return $CY\operatorname{sgn}(\gamma_S)\chi_S(X)$. The empty term uses no
feature; if $C=0$, return zero. This yields the required observable score and bound. The
finite-sample radius can therefore be minimized by a linear program:
\[
 \min_{\gamma,\varepsilon\ge0}
 \varepsilon+\sqrt{2\log(2/\alpha)/n}\|\gamma\|_1,
 \qquad |d-\mathsf F\gamma|\le\varepsilon\mathbf1,
\]
where only $\varepsilon$ is constrained nonnegative and $\gamma$ is free.
This is a conservative confidence criterion, distinct from the MSE
criterion in Appendix~\ref{app:dense}.

\paragraph{What scales without full form enumeration.}
The all-form search certifies an offline benchmark; deployment of an
already fitted score requires no such search. If a frozen contrast has
a supplied decomposition $d=\sum_j w_j f_j$, with each $f_j$ computable
within $b$ queries and $|f_j|\le M_j$, choose $j$ with probability
$|w_j|M_j/C$ for $C=\sum_j|w_j|M_j$ and return
$CY\operatorname{sgn}(w_j)f_j/M_j$.
It is exactly unbiased and bounded by $C$. This includes differences of
probability-averaging forests with depth at most $b$, even when their
combined prediction uses many coordinates and the inputs are continuous.
A fixed form can simulate a sampled component by its path-set construction.
Neither method asserts optimal efficiency. Nonlinear transformations of
ensemble scores, including the logistic link in ordinary GBDT probability
outputs, need a separate approximation certificate.

For an arbitrary continuous model, restrict selected coordinates to two
values and hold the others fixed. The restriction of any $b$-coordinate
observable function has Boolean degree at most $b$. Thus a bias-floor
lower certificate on that finite slice is also a lower bound on the full
product-domain problem. The slice supplies no full-domain upper bound.
This is a one-sided rejection tool whose scope must be stated explicitly.

\subsection{Independent GBDT and Neural-Network Versions}\label{app:nonlinear}
WDBC contains 569 diagnostic records and 30 image-derived continuous
features \citep{wolberg1993}. The target is malignant versus benign.
The data therefore supply a non-questionnaire outcome and continuous
model inputs, but their image-derived acquisition does not instantiate
our survey cap. This experiment asks about mathematical feasibility for
ordinary model classes.

Five fixed seeds, 20261108--20261112, permute the data. The first 100
records train the old model; the next 300 train the new model; the last
169 determine its evaluated accuracy gain. GBDT uses 30 versus 100 trees,
depth two, learning rate 0.05 and minimum leaf size five. The MLP uses
one ReLU hidden layer of width eight versus sixteen, $\ell_2$ penalty one,
and L-BFGS. Each model's standardization uses its own training records.
No evaluation record selects model size, variables, thresholds or seed. The splits reuse individuals and are descriptive sensitivity analyses, not independent populations.

The discrete condition takes twelve coordinates at equally spaced indices
of the original feature list and thresholds them at medians of the 400
training records. Both probability outputs and deterministic thresholds
at one half are evaluated. The admissible domain is the entire binary
cube, with $b=3$. The continuous condition trains on all thirty original
features. Its certificate varies the same twelve coordinates between
their training first and third quartiles and fixes the others to training
medians. Continuous-condition brackets refer to that slice; only their
lower endpoints apply to the full domain. Arbitrary product combinations
can violate empirical feature dependencies, so these are distribution-free
certificates under the declared domain, not support-aware clinical claims.

For a target vector $d$ on the 4,096-profile grid and the degree-three
Walsh matrix $\mathsf F$, the floor program minimizes
$\|d-\mathsf F\gamma\|_\infty$. Its floating-point solution is converted
to certificates: round primal coefficients to multiples of $2^{-28}$;
round a dual vector to integers $z$ on scale $2^{40}$; then form
\[
 h_{\rm int}=4096z-\mathsf F\mathsf F^Tz.
\]
Because $\mathsf F^T\mathsf F=4096I$, the integer vector has exactly zero
degree-at-most-three coefficients. Dividing by its exact $\ell_1$ norm
gives a feasible rational dual. Both its inner product with the stored
binary64 predictions and the rounded primal's maximum residual are
computed with rational arithmetic. The certificate treats those stored
software predictions as the defined finite target, without claiming
formal verification of a real-arithmetic neural network.

The binary probability condition also minimizes the bias-expanded
Hoeffding radius at $n=1000$, $\alpha=.05$. The exact certificates do not
depend on an optimizer's success flag. A bounded-time implementation can
fall back to a feasible low-degree projection and its orthogonal residual;
solver status and certificate width are released in every row.
All thirty floor programs terminated optimally in the released run. The
largest certified bracket width is below $10^{-7}$, and the fitted models
emitted no convergence warnings.

\begin{table*}[t]
\centering
\caption{Cap-three nonlinear-model diagnostic. Binary rows bracket the full binary-domain floor. Continuous rows bracket a twelve-coordinate slice; only the lower endpoints transfer to the full continuous domain. The last column compares the certified floor directly with the requested half-width.}
\label{tab:nonlineardetail}
\small
\setlength{\tabcolsep}{4pt}
\begin{tabular}{lrrrr}
\toprule
Model condition & Splits & Median LB & Median UB & LB$>0.025$ \\
\midrule
Binary, probability, GBDT & 5 & 0.0552 & 0.0552 & 5/5 \\
Binary, probability, MLP & 5 & 0.2780 & 0.2780 & 5/5 \\
Binary, threshold, GBDT & 5 & 0.3409 & 0.3409 & 5/5 \\
Binary, threshold, MLP & 5 & 0.6572 & 0.6572 & 5/5 \\
Continuous slice, GBDT & 5 & 0.0065 & 0.0065 & 0/5 \\
Continuous slice, MLP & 5 & 0.1286 & 0.1286 & 5/5 \\
\bottomrule
\end{tabular}
\end{table*}


The continuous slice results are weaker statements than the complete
binary-cube results. A small slice lower bound cannot certify an acceptable
full-domain bias. Similarly, reporting a good empirical MSE does not
replace the uniform residual required for a bias-expanded interval.
